\documentclass[11pt,a4paper]{article}
\usepackage[margin=2.3cm]{geometry}
\usepackage{amsmath,amssymb}
\usepackage{xcolor}
\usepackage{xr}
\usepackage[hidelinks]{hyperref}
\externaldocument{main}
\setlength{\parindent}{0pt}\setlength{\parskip}{5pt}
\newcommand{\cmt}[1]{\par\noindent{\small\itshape Editor's comment: #1}\par}
\newcommand{\resp}[1]{\par\noindent\textbf{Response.} #1\par}
\newcommand{\loc}[1]{\par\noindent\textbf{Changes.} #1\par}
\begin{document}
\begin{center}
{\Large\bfseries Response to the editor and reviewers (second revision)}\\[4pt]
Manuscript ID AISBE-08-2026-0041.R1\\
\emph{Physics-informed e-processes for damage detection under a lifetime false-alarm budget in continuously monitored structures}
\end{center}

Dear Dr. Khudhair,

Thank you for the editorial assessment and for the reviewers' recommendations. Each point of your letter is answered below, in the order in which it was raised, with the location of the change in the revised manuscript. Section, table, figure and proposition numbers refer to the revised manuscript. In the highlighted version, text added or changed in this round is set in red; text from the previous round remains blue; the clean version carries no colour. Several of your points concerned numbers, so we re-derived them from the stored results and, where a number depended on a simulation stage, reran that stage with its fixed seeds; what we found is stated under each point.

Summary. All points were addressed by changes to wording, proofs, tables and supplied files. One numerical statement was wrong in the manuscript (the observation window of the stiffening-then-damage experiment); the others were reporting and labelling errors. No headline result changed. The conformal comparators are now described as empirical comparators without a sequential guarantee, as you allowed, so no simulation had to be changed for that point.

\section*{Editorial points}

\subsection*{Editorial point 1. Conformal comparators (Section 4.1.3, Table 4)}

\cmt{Repeatedly ranking monitoring scores against the same finite calibration set does not, under exchangeability or i.i.d. sampling alone, give independent uniform p-values; the argument does not establish a test-martingale guarantee for the product of the power betting functions. Provide the p-value construction and a valid justification, or describe CTM and CTM-F as empirical comparators; revise Table 4 and the interpretation of their false-alarm behaviour.}

\resp{We agree, and we have chosen the second option. The implemented construction is now stated exactly: with calibration scores $c_1,\dots,c_n$ fixed at the end of commissioning, the smoothed conformal p-value is $p_t=\dfrac{\#\{i:c_i>X_t\}+U_t\,(1+\#\{i:c_i=X_t\})}{n+1}$ with independent uniform U\_t, floored at $10^{-6}$; the betting functions $\epsilon\,p^{\epsilon-1}$ are multiplied over epochs for twelve equally spaced $\epsilon\in[0.05,0.95]$, averaged with equal weights, and the alarm is at $1/\alpha$. This is what the code (monitors.py) computes.}

\par We then state why no guarantee is claimed. Conditional on the realised calibration set, the distribution of p\_t is the empirical distribution function of that set evaluated at the new score, which is in general not uniform (deviations of order $n^{-1/2}$). Marginally over the calibration set, p-values of different epochs share it and are dependent. The conditional expectation of the betting function need not be one, so the product is not a martingale and Ville's inequality does not apply. A construction with a sequential guarantee would rank each score among all earlier scores, calibration scores included (online conformal p-values), and tests exchangeability of the whole sequence, a different null; we did not implement it.

\par The earlier sentence that the p-values ``are uniform and independent, and the product is a test martingale'' if the scores are exchangeable with a fixed calibration set has been removed, and the description of CTM and CTM-F as ``anytime-valid comparators'' has been removed everywhere. In Table~\ref{tab:fwer} they now form a separate group, ``conformal comparators: empirical, no established anytime-valid guarantee'', and the caption says why. The text of Section~\ref{sec:res-fa} no longer calls the excess of CTM (0.033 and 0.030) a failure of a guarantee, and no longer reads the zero of CTM-F as certification; the comparison with PHASE now rests on the fact that PHASE has a proved conditional guarantee and the conformal comparators do not. The discussion of the power frontier in Section~\ref{sec:res-delay} was adjusted in the same way. The simulation results for CTM and CTM-F are unchanged, because the implementation is unchanged.

\loc{Section~\ref{sec:comparators} (rewritten, the formula for p\_t and the explanation are new); Table~\ref{tab:fwer} (new group and caption); Section~\ref{sec:res-fa} (paragraph on the conformal comparators); Section~\ref{sec:res-delay} (frontier paragraph); Section~\ref{sec:related-seq} and Table~\ref{tab:criteria} (the guarantee is attributed to the online construction only); Contribution 5 in Section~\ref{sec:intro}; Appendix~\ref{app:notation} and Appendix~\ref{app:impl} (formula and status of CTM and CTM-F).}

\subsection*{Editorial point 2. Propositions 2, 4(iii) and 5 (Appendix C)}

\cmt{Proposition 2: the lower bound is stated for the raw maximum but the algorithm uses its positive part (counterexample: negative real ray, $V=\{-1\}$, r = 1); state the result for the positive part, treat the zero-projection case separately. Proposition 4(iii): state $0<\varepsilon<\alpha$. Proposition 5: make the auxiliary hitting time consistent with the case in which it equals the change time.}

\resp{Proposition 2. We checked your counterexample: the cone is $\{x\le0\}$, $V=\{-1\}$, covering radius zero, r = 1 gives a raw maximum of -1 and a projection of norm 0, so the lower bound as stated was false for the raw maximum. The proposition is now stated for $\big(\max_k\langle v_k,r\rangle\big)^{+}$, which is what the score (eq. 9) uses; the upper bound holds for both. The proof treats $p=\|\Pi_{\mathcal C}(r)\|=0$ separately (the bound is then trivial because the positive part is non-negative) and, for p $>$ 0, uses the unit vector $\Pi_{\mathcal C}(r)/p$, which also removes the earlier tacit assumption that the supremum is attained on the unit sphere. The counterexample is recorded in the proof.}

\par Proposition 4(iii). The condition is now $0<\varepsilon<\alpha$, so that the threshold $1/(\alpha-\varepsilon)$ is finite and exceeds $1/\alpha$. The experiments use $\varepsilon=\alpha/2$, which satisfies it.

\par Proposition 5. The auxiliary hitting time is now defined as $T=\inf\{t\ge t^\star:L^\lambda_t\ge b\}$, so that T = t* when the component already satisfies $L^\lambda_{t^\star}\ge b$; the proof treats that case first (the right side of the bound is non-negative, so it holds trivially) and otherwise has T $>$ t*, with the log-wealth below b before T. The first-passage argument is otherwise unchanged.

\loc{Proposition~\ref{prop:poly} and its proof in Appendix~\ref{app:proofs}; Proposition~\ref{prop:delta}(iii) and its proof; Proposition~\ref{prop:delay} proof.}

\subsection*{Editorial point 3. Consistency of the qualifications on the false-alarm guarantee}

\cmt{The abstract makes coverage of the operating environment sufficient for the lifetime bound; the last sentence of Section 3.3 implies that predictable fitting alone preserves validity; the statement in Section 3.5 about acting before threshold crossing should agree with Section 6.1. Retain these distinctions when explaining the PHASE acronym and summarising performance.}

\resp{Abstract. The practical-implications field now states the necessary conditions: the bound holds only if the expected score of the undamaged structure stays below the design constant plus the calibrated allowance, which requires commissioning that represents the whole operating environment; coverage is necessary but does not by itself establish the bound or its calibration; missed detections are not controlled.}

\par Section 3.3. The last sentence now says that predictability of the plug-in is necessary but not sufficient: validity also needs the conditional-mean bound for the score formed from the residual, which a predictable fit does not deliver and which Section 3.6 prices.

\par Section 3.5. The sentence now says that the bound holds whenever and however often the operator inspects the process, concerns the alarm at level $1/\alpha$, and that an action at a lower evidence level $w$ is protected only at level $1/w$; it points to Section 6.1, which already said so.

\par PHASE acronym and summaries. Where the acronym is introduced (abstract and Section 1) the text now says that anytime-validity is a conditional property of the e-process and that, under the empirical calibration of the allowance, the observed false-alarm control is an empirical finding and not an instance of the theorem; the same reservation is given for the title's ``lifetime false-alarm budget''. This also answers Reviewer 1's second proof-stage item. We also added the qualification that coverage is necessary but not sufficient in Section 6.1 (third implication) and in the conclusions.

\loc{Abstract (Design/methodology/approach and Practical implications); Section~\ref{sec:intro} (paragraph introducing PHASE); Section 3.3 (final sentences); Section~\ref{sec:eprocess} (second feature); Section~\ref{sec:practice}; Section~\ref{sec:deploy} (``Are the delays acceptable?''); Section~\ref{sec:concl}.}

\subsection*{Editorial point 4a. Table 3, error measure}

\cmt{Table 3 is described as the maximum relative error across modes, but an independent calculation reproduces the values as maximum absolute approximation error divided by maximum absolute exact shift. For 10\% loss at storey 3 these are 8.23\% (the table) and 13.87\% (maximum of the modal relative errors). Define the measure and align caption and interpretation.}

\resp{You are right. We recomputed from the stored exact and first-order shifts: the table reported the normalised error $\max_i|f_i-e_i|/\max_i|e_i|$ (8.23\% for 10\% loss at storey 3) and the caption called it a relative error; the maximum of the six modal relative errors is 13.87\% (sixth mode). Both measures are now defined in the text and the caption. Table 3 has two panels: A, the normalised error (values unchanged: 1.9, 1.7, 1.8\% at 2\% loss; 4.6, 4.2, 4.6\% at 5\%; 9.3, 8.2, 9.2\% at 10\%) and B, the maximum modal relative error (1.9, 2.8, 2.4\%; 4.6, 7.0, 5.8\%; 9.3, 13.9, 11.6\%).}

\par The interpretation was changed to match. The statement that the first-order model is ``accurate'' with errors ``degrading gracefully and conservatively'' has been replaced: the model is adequate at 2-5\% loss (individual modal errors of a few per cent) but its error at 10\% loss is not negligible (up to 13.9\% of an individual modal shift); the first-order model understates the frequency drop in most mode-storey pairs; the word ``conservatively'' was removed because we cannot show that the direction of the error is conservative for detection. The code now stores both measures (stage.py) and table 3 is generated by make\_tables.py.

\loc{Section~\ref{sec:sim-results} (subsection ``Validity of the first-order physics''): text and Table~\ref{tab:validation}.}

\subsection*{Editorial point 4b. Observation windows in Section 4.2.9 and Table 12}

\cmt{Detection is stated to be counted within nine months, while Table 12 gives control-case median delays of 300 days (PHASE-C) and 368 days (PHASE-Omega). Reconcile the windows, detection counts and delay summaries and make differences between control and stiffening cases explicit.}

\resp{The inconsistency was real. We reran the stage with its fixed seeds (it reproduces every entry of the old table) and recomputed detection in both windows. The experiment uses four-year records (5,840 epochs) and the loss occurs at epoch 2,545, so 3,295 epochs (824 days, about 27 months) remain after the loss; the original counts were over that whole remaining record and not over nine months as the text said. That is why medians of 300 and 368 days are possible. The nine-month window of the other experiments is 1,095 epochs.}

\par Table 12 now reports both windows for the control (loss alone). Within nine months: PHASE 100/100 (median 92 days), PHASE-R 98/100 (12), PHASE-C 15/100 (254), CTM-F 90/100 (220), CUSUM 100/100 (24), PHASE-Omega 3/100 (249). Over the whole remaining record: PHASE 100/100 (92), PHASE-R 98/100 (12), PHASE-C 70/100 (300), CTM-F 100/100 (226), CUSUM 100/100 (24), PHASE-Omega 14/100 (368). For stiffening followed by loss the result does not depend on the window: no monitor detected the loss, and only PHASE-Omega alarmed, on the stiffening itself, in 100/100 records at a median of 11 days (not counted as a detection of the loss, by the definition in Section 4.1.4). The text states the window, explains why the medians exceed nine months, and states that the control and the stiffening cases differ only in what precedes the loss. The definition of detection in Section 4.1.4 notes this experiment as the exception to the nine-month window, and Section 6.2 now says ``within nine months or over the 27 months that remained''.

\par The new script stiffdam\_windows.py produces the two-window counts from the stage output; make\_tables.py generates the table.

\loc{Section~\ref{sec:res-masked} and Table~\ref{tab:masked}; Section~\ref{sec:metrics} (detection definition); Section~\ref{sec:deploy} (long-lived assets); code/stiffdam\_windows.py, code/stiffdam\_windows.json.}

\subsection*{Editorial point 4c. Zero-event confidence limits in Table 9}

\cmt{The two-sided exact 95\% upper limit for 0/100 is about 0.0362, not 0.03.}

\resp{Confirmed. The stored limit (summary.json) was 0.0362 and the caption of Table 8 already said 0.036; the table generator had truncated limits to two decimals instead of rounding them. Table 9 now shows [0.00, 0.04] for every 0/100 entry. We checked every k/n [lo, hi] interval in the manuscript against an exact (Clopper-Pearson) computation and corrected the other truncated limits: Table 5 (150/150: lower limit 0.9757, now 0.98, was 0.97; the same in the text of Section 4.2.3) and Table 9 (66/100: lower limit 0.56, was 0.55; 99/100: 0.95, was 0.94; 92/100: 0.85, was 0.84). The generator was corrected (make\_tables.py) and regenerates the tables identically to the manuscript. All other intervals were already correctly rounded.}

\loc{Table~\ref{tab:late}; Table~\ref{tab:delaydetail} and Section~\ref{sec:res-delay} (text); code/make\_tables.py.}

\subsection*{Editorial point 5. Field results}

\cmt{Restrict the Z24 statement that a frequency rise cannot have been caused by damage to the assumed stiffness-loss mechanism at fixed mass; keep the KW51 pre-works condition described as assumed, and the configuration counts descriptive, not independent estimates of false-alarm probability.}

\resp{Z24. The sentence in the directional-behaviour paragraph is restricted: under the model on which the cone rests (stiffness loss at fixed mass) a frequency rise cannot have been caused by damage; it does not follow that no damage of any kind occurred, since mass loss, crack closure or boundary stiffening can raise frequencies and lie outside the cone, as Section 5.2 already recognised.}

\par KW51. The pre-works condition is described as an assumed healthy reference throughout; an alarm before the inferred works period is called a ``false alarm'' only relative to that reference (and would not be false if the structure was not healthy). The text, the Table 17 caption, the abstract, Table 14, the real-data summary, the limitations and the conclusions use this wording. The counts over the 15 configurations are stated to be descriptive: the configurations share one bridge and one pre-works record, overlap, and were selected exploratorily, so they are not independent estimates of a false-alarm probability.

\loc{Section~\ref{sec:z24-results} (directional behaviour); Section~\ref{sec:kw51} (setup, results, ``What this shows''; Table~\ref{tab:kw51} caption); Table~\ref{tab:components}; Section~\ref{sec:z24-adds}; Section~\ref{sec:limits}; Section~\ref{sec:concl}; Abstract.}

\subsection*{Editorial point 6. Implementation files and reproduction instructions}

\cmt{Make the files cited in Appendix B and the response accessible through submission files or a documented repository.}

\resp{The code is supplied in two ways. (i) As Supplementary file S1 (phase\_code.zip) with this submission: the code directory, the extracted Z24 modal features, the KW51 results file, the vector source of Figure 10 and README\_reproduction.md, which lists every stage and command. (ii) In a documented repository, https://github.com/vinhqdang/Physics-informed-e-processes-for-anytime-valid-damage-detection-in-continuously-monitored-structures (branch main), which holds the manuscript sources, figures, code, the previous and present responses, and the same README. Appendix B now has an ``Availability'' paragraph naming the archive and stating that the repository address is withheld from the anonymised manuscript only and will be cited in the published version. The raw Z24 records (KU Leuven) and the KW51 tracked-mode file (Zenodo 3745914, CC BY-NC-SA 4.0) are not redistributed; the README gives provenance and the download command.}

\par You could not locate these files in the submission's file list; the archive and the repository are the correction of that.

\loc{Appendix~\ref{app:impl} (``Availability''); README.md, README\_reproduction.md, supplementary/phase\_code.zip.}

\subsection*{Editorial point 7. Section 6.2, ``Are the delays acceptable?''}

\cmt{The paragraph is incomplete and should be restored.}

\resp{The paragraph was truncated by a typesetting error: the percent signs and a cross-reference in it had been escaped, so the LaTeX compiler treated the text after ``5'' in ``5\% loss'' as a comment and dropped the rest. It is restored in full in both the highlighted and clean versions. We checked the whole manuscript for the same fault and found no other instance. We also added that the reported delays are the price of targeting a 1\% lifetime budget that holds only under the conditions of Section 3.6.}

\loc{Section~\ref{sec:deploy}, paragraph ``Are the delays acceptable?''.}

\subsection*{Editorial point 8. Title, Figure 10 and placement of the appendices}

\cmt{Synchronise the revised title across the manuscript and submission metadata; provide a suitable source version of Figure 10 with labels legible at the final size; agree the placement of the appendices with the editorial office.}

\resp{Title. The revised title (``Physics-informed e-processes for damage detection under a lifetime false-alarm budget in continuously monitored structures'') is the title in the manuscript, its PDF metadata, the response letters and the supplied files. The repository name keeps the first-submission title and cannot be changed without breaking the address; the README says so. SUBMISSION\_METADATA.txt gives the title, keywords and file list to be entered in the submission system. Because we cannot edit the system's metadata from here, please confirm that the entry shows the revised title; we will also check it at proof stage, as Reviewer 1 asked.}

\par Figure 10 (KW51). It is now supplied as a vector PDF (fig\_kw51.pdf, TrueType fonts embedded), 7.0 in (17.8 cm) wide, the width of a two-column page, with all text 8 pt or larger at that size (legend entries were 6 pt before and now sit outside the plotting area, so no label overlaps data). The script that draws it (kw51\_figure.py) is supplied; a PNG preview is kept. The other figures are unchanged raster images and we can supply vector sources if the production team asks.

\par Appendices. The manuscript places Appendices A (notation and acronyms), B (implementation details) and C (proofs) after the reference list, at the end of the file. We have not changed the layout. Please tell us whether the journal wishes them printed as appendices, moved to the main text (the proofs are referred to throughout), or supplied as online supplementary material; we will follow the editorial office's decision.

\loc{Title page and PDF metadata; Figure~\ref{fig:kw51} (figures/fig\_kw51.pdf); appendices at the end of the manuscript; SUBMISSION\_METADATA.txt.}

\section*{Reviewers}

Reviewer 1. We thank the reviewer. The three proof-stage items are handled as follows. (1) Title in the publisher's metadata: see editorial point 8; we will verify it at proof. (2) Clarification where the PHASE acronym is introduced about anytime-validity under the empirical calibration: added in the abstract and in Section 1 (editorial point 3). (3) Figure 10 in vector format with legible text at two-column size: supplied (editorial point 8).

Reviewer 2. We thank the reviewer; no further comments were raised.

We hope the revised manuscript can now be accepted. We are grateful to the editor and the reviewers for their care, and in particular for the checks of the proofs and numbers that led to these corrections.

\end{document}
